\documentclass[12pt]{article}

\usepackage[a4paper,margin=1in]{geometry}
\usepackage{amsmath}
\usepackage{enumitem}

\begin{document}

\begin{center}
    {\LARGE Circular Motion Worksheet}\\[0.5em]
\end{center}

\vspace{1em}

\begin{enumerate}[label=\textbf{\arabic*.}, itemsep=2em]

\item
A \(0.40\,\text{kg}\) ball is attached to a string and moves in a horizontal circle of radius \(1.5\,\text{m}\) on a frictionless table. The ball completes \(2.5\) revolutions every second.

\begin{enumerate}[label=(\alph*)]
    \item Calculate the angular velocity \(\omega\) of the ball.
    \item Determine the speed of the ball.
    \item Find the tension in the string.
\end{enumerate}

\item
A small mass of \(0.75\,\text{kg}\) is attached to a string and whirled in a horizontal circle of radius \(0.80\,\text{m}\). The tension in the string is \(48\,\text{N}\).

\begin{enumerate}[label=(\alph*)]
    \item Calculate the centripetal acceleration of the mass.
    \item Determine the angular velocity \(\omega\).
    \item Find the time period of the motion.
\end{enumerate}

\item
A car of mass \(1200\,\text{kg}\) travels around a flat circular track of radius \(55\,\text{m}\) at a constant speed of \(20\,\text{m s}^{-1}\).

\begin{enumerate}[label=(\alph*)]
    \item Calculate the centripetal acceleration of the car.
    \item Determine the resultant centripetal force required to keep the car moving in a circle.
    \item State which force provides the centripetal force in this situation.
\end{enumerate}

\item
A \(0.20\,\text{kg}\) mass attached to a light string moves in a vertical circle of radius \(0.90\,\text{m}\) at a constant speed of \(5.0\,\text{m s}^{-1}\).

\begin{enumerate}[label=(\alph*)]
    \item Calculate the required centripetal force.
    \item Determine the tension in the string at the bottom of the circle.
    \item Determine the tension in the string at the top of the circle.
\end{enumerate}

\item
A bucket of water of total mass \(1.8\,\text{kg}\) is swung in a vertical circle of radius \(1.2\,\text{m}\). At the top of the circle the tension in the handle is \(14\,\text{N}\).

\begin{enumerate}[label=(\alph*)]
    \item Draw a force diagram for the bucket at the top of the circle.
    \item Calculate the resultant force toward the centre at the top of the circle.
    \item Determine the speed of the bucket.
\end{enumerate}

\newpage
\item
A rollercoaster car of mass \(450\,\text{kg}\) moves around a vertical circular track of radius \(18\,\text{m}\). At the bottom of the loop the normal contact force exerted by the track on the car is \(1.25\times10^4\,\text{N}\).

\begin{enumerate}[label=(\alph*)]
    \item Calculate the weight of the car.
    \item Determine the centripetal force acting on the car at the bottom of the loop.
    \item Calculate the speed of the car.
\end{enumerate}

\item
A car travels around a banked curve of radius \(85\,\text{m}\). The banking angle is \(12^\circ\). Assume there is no friction.

\begin{enumerate}[label=(\alph*)]
    \item Draw a force diagram for the car.
    \item Resolve the normal reaction force vertically and horizontally.
    \item Show that the speed required to negotiate the bend without friction is
    \[
    v = \sqrt{rg\tan\theta}
    \]
    \item Calculate the required speed.
\end{enumerate}

\item
A motorcycle travels around a frictionless banked curve of radius \(40\,\text{m}\) at a speed of \(16\,\text{m s}^{-1}\).

\begin{enumerate}[label=(\alph*)]
    \item Determine the required centripetal acceleration.
    \item Calculate the banking angle needed for the turn.
    \item State why friction is not required at this speed.
\end{enumerate}

\item
A \(1500\,\text{kg}\) car travels around a frictionless banked bend of radius \(60\,\text{m}\). The banking angle is \(9^\circ\).

\begin{enumerate}[label=(\alph*)]
    \item Calculate the normal reaction force acting on the car.
    \item Determine the horizontal component of the normal force.
    \item Calculate the speed of the car.
\end{enumerate}

\end{enumerate}

\newpage

\begin{center}
    {\LARGE Answer Key}
\end{center}

\vspace{1em}

\begin{enumerate}[label=\textbf{\arabic*.}, itemsep=2em]

\item
\begin{enumerate}[label=(\alph*)]
    \item
    \[
    \omega = 2\pi f = 2\pi(2.5)=15.7\,\text{rad s}^{-1}
    \]

    \item
    \[
    v=r\omega = 1.5(15.7)=23.6\,\text{m s}^{-1}
    \]

    \item
    \[
    T=\frac{mv^2}{r}
    =\frac{0.40(23.6)^2}{1.5}
    \approx 148\,\text{N}
    \]
\end{enumerate}

\item
\begin{enumerate}[label=(\alph*)]
    \item
    \[
    F_c=ma_c
    \]
    \[
    a_c=\frac{48}{0.75}=64\,\text{m s}^{-2}
    \]

    \item
    \[
    a_c=r\omega^2
    \]
    \[
    \omega=\sqrt{\frac{64}{0.80}}
    =8.94\,\text{rad s}^{-1}
    \]

    \item
    \[
    T=\frac{2\pi}{\omega}
    =0.703\,\text{s}
    \]
\end{enumerate}

\item
\begin{enumerate}[label=(\alph*)]
    \item
    \[
    a_c=\frac{v^2}{r}
    =\frac{20^2}{55}
    =7.27\,\text{m s}^{-2}
    \]

    \item
    \[
    F_c=ma_c
    =(1200)(7.27)
    \approx 8.73\times10^3\,\text{N}
    \]

    \item Friction between the tyres and the road provides the centripetal force.
\end{enumerate}

\item
\begin{enumerate}[label=(\alph*)]
    \item
    \[
    F_c=\frac{mv^2}{r}
    =\frac{0.20(5.0)^2}{0.90}
    =5.56\,\text{N}
    \]

    \item Bottom:
    \[
    T-mg=F_c
    \]
    \[
    T=5.56+1.96
    =7.52\,\text{N}
    \]

    \item Top:
    \[
    T+mg=F_c
    \]
    \[
    T=5.56-1.96
    =3.60\,\text{N}
    \]
\end{enumerate}

\item
\begin{enumerate}[label=(\alph*)]
    \item Both tension and weight act toward the centre.

    \item
    \[
    F_c=T+mg
    \]
    \[
    F_c=14+(1.8)(9.81)
    =31.7\,\text{N}
    \]

    \item
    \[
    \frac{mv^2}{r}=31.7
    \]
    \[
    v=\sqrt{\frac{31.7(1.2)}{1.8}}
    \approx 4.60\,\text{m s}^{-1}
    \]
\end{enumerate}

\item
\begin{enumerate}[label=(\alph*)]
    \item
    \[
    W=mg=(450)(9.81)=4410\,\text{N}
    \]

    \item
    \[
    F_c=R-W
    \]
    \[
    F_c=12500-4410
    =8090\,\text{N}
    \]

    \item
    \[
    \frac{mv^2}{r}=8090
    \]
    \[
    v=\sqrt{\frac{8090(18)}{450}}
    \approx 18.0\,\text{m s}^{-1}
    \]
\end{enumerate}

\item
\begin{enumerate}[label=(\alph*)]
    \item ---

    \item
    \[
    R\cos\theta = mg
    \]
    \[
    R\sin\theta = \frac{mv^2}{r}
    \]

    \item Dividing equations:
    \[
    \tan\theta=\frac{v^2}{rg}
    \]
    therefore
    \[
    v=\sqrt{rg\tan\theta}
    \]

    \item
    \[
    v=\sqrt{(85)(9.81)\tan12^\circ}
    \approx 13.3\,\text{m s}^{-1}
    \]
\end{enumerate}

\item
\begin{enumerate}[label=(\alph*)]
    \item
    \[
    a_c=\frac{v^2}{r}
    =\frac{16^2}{40}
    =6.4\,\text{m s}^{-2}
    \]

    \item
    \[
    \tan\theta=\frac{v^2}{rg}
    \]
    \[
    \theta=\tan^{-1}\left(\frac{16^2}{40(9.81)}\right)
    \approx 33.1^\circ
    \]

    \item The horizontal component of the normal reaction alone provides the required centripetal force.
\end{enumerate}

\item
\begin{enumerate}[label=(\alph*)]
    \item
    \[
    R\cos9^\circ = mg
    \]
    \[
    R=\frac{1500(9.81)}{\cos9^\circ}
    \approx 1.49\times10^4\,\text{N}
    \]

    \item
    \[
    R\sin9^\circ
    \approx 2330\,\text{N}
    \]

    \item
    \[
    \frac{mv^2}{r}=2330
    \]
    \[
    v=\sqrt{\frac{2330(60)}{1500}}
    \approx 9.66\,\text{m s}^{-1}
    \]
\end{enumerate}

\end{enumerate}

\end{document}